\Needspace{12\baselineskip}
\begin{center}
\textit{Historical remarks.}
\end{center}

\Needspace{8\baselineskip}
\Article{222}
Since several of the matters we have discussed so far have also been considered by other mathematicians, we cannot pass over their contributions in silence. The distinguished La Grange made general investigations into the \textit{equivalence of forms}, \textit{Nouv. M\'em. de l'Ac. de Berlin,} 1773 \textit{p.} 263 and 1775 \textit{p.} 323 ff. In particular, he showed that for any given determinant there is a finite number of forms such that every form with that determinant is equivalent to one of them, and so all forms with a given determinant can be divided into classes. Later, the celebrated Le Gendre discovered several elegant properties of this classification, mostly by induction from observed cases; we will give these below and support them with proofs. But no one had yet reached the distinction between proper and improper equivalence, whose usefulness is most apparent in the finer investigations.

The distinguished La Grange was the first to solve completely\SourceLeaf{227} the famous problem explained in Article 216 ff., \textit{Hist. de l'Ac. de Berlin,} 1767 \textit{p.} 165 and 1768 \textit{p.} 181 ff. A solution (but a less complete one) also appears in the \textit{Suppl. ad Euleri Algebram}, already cited several times. Earlier, the distinguished Euler had tackled the same subject, \textit{Comm. Petr. T.} VI \textit{p.} 175; \textit{Comm. Nov. T.} IX \textit{p.} 3; \textit{Ibid. T.} XVIII \textit{p.} 185 ff. But he always restricted his investigation to deriving other solutions from some solution he assumes already known; moreover, his methods can provide \textit{all} solutions in only a few cases (see La Grange, \textit{Hist. de l'Ac. de Berlin} 1767, \textit{p.} 237). Since the last of these three papers is later than La Grange's solution, which covers the problem in full generality and leaves nothing to be desired in this respect, Euler seems not yet to have known that solution at the time (Volume XVIII of the Commentaries belongs to the year 1773 and was published in 1774). Our solution, however (like everything else we have given so far in this section), rests on entirely different principles.

The relevant results given by others, Diophantus, Fermat, etc., concern very special cases. Since we have already mentioned above those that seemed particularly worth recording, we will refrain from listing everything individually.

\begin{center}
$\begin{array}{ccc}\ast&&\ast\\&\ast&\end{array}$
\end{center}

What we have set out so far about forms of the second degree should be regarded only as the first elements of this theory. As we pursued this investigation more closely, a vast field opened up before us; in what follows we will select from it those matters that seem especially worth attention. For this subject is so fertile that, for brevity, we must pass over in silence many other things we have already managed to discover; far more, without doubt, still lie hidden and await new efforts. At the start of these investigations, we should note at once that forms with determinant $0$ are excluded unless we state otherwise.

\Needspace{12\baselineskip}
\begin{center}
\textit{FURTHER INVESTIGATIONS INTO FORMS.}
\end{center}

\Needspace{12\baselineskip}
\begin{center}
\textit{Division of forms with a given determinant into classes.}
\end{center}

\Needspace{8\baselineskip}
\Article{223}
We have already shown above (Articles 175, 195, 211) that for any given integer $D$ (positive or negative), we can specify a finite number of forms $F,F',F''$ etc. with determinant $D$, such that every form with determinant\SourceLeaf{228} $D$ is properly equivalent to one, and only one, of them. Thus all forms with determinant $D$ (of which there are infinitely many) can be \textit{classified} according to those forms: the collection of all forms properly equivalent to $F$ makes up the first class; those properly equivalent to $F'$ make up the second, etc.

From each class of forms with a given determinant $D$, we can choose a form and regard it as a \textit{representative form} of the whole class. In itself, the choice of form from each class is entirely arbitrary, but we should always prefer one that seems to surpass the others in \textit{simplicity}. The simplicity of a form $(a,b,c)$ must clearly be judged from the size of the numbers $a,b,c$, and we may rightly call a form $(a',b',c')$ less simple than $(a,b,c)$ if $a'>a$, $b'>b$, $c'>c$. But this does not settle the matter completely: for \textit{example}, it remains our choice which of the forms $(17,0,-45)$, $(5,0,-153)$ we prefer to call simpler. Usually, however, it will be useful to follow this rule:

I. When the determinant $D$ is negative, take the reduced forms contained in the individual classes as their representative forms. If two reduced forms occur in the same class (which will be opposite, Article 172), take the one whose middle term is positive.

II. When the determinant $D$ is positive and nonsquare, work out the period of a reduced form contained in the proposed class. This period will contain either two ambiguous forms or none (Article 187).

1) In the first case, let the ambiguous forms be $(A,B,C)$, $(A',B',C')$; let the least residues of the numbers $B,B'$ modulo $A,A'$ be $M,M'$, respectively (these can be taken positively unless they are $=0$); finally, let $\frac{D-MM}{A}=N$, $\frac{D-M'M'}{A'}=N'$. With this done, take whichever of the forms $(A,M,-N)$, $(A',M',-N')$ seems simplest as the representative form. In deciding this, prefer a form whose middle term is $=0$. If the middle term is $0$ in both or in neither, prefer the form with the smaller first term; and when the first terms are equal in size but opposite in sign, prefer the positive sign to the negative.

2) If there is no ambiguous form in the whole period, choose from all the forms of the period the one with the smallest first term, disregarding its sign\SourceLeaf{229}. If two forms occur in the same period with the same first term positive in one and negative in the other, prefer the former to the latter. Let this form be $(A,B,C)$, and derive from it, in the same way as in the preceding case, another form $(A,M,-N)$ (namely, by taking $M$ to be the absolutely least residue of $B$ modulo $A$, and setting $N=\frac{D-MM}{A}$). Take this last form as the representative.

If the same smallest first term $A$ happens to be shared by several forms of the period, treat all these forms in the way prescribed, and take as representative the resulting form whose middle term is as small as possible.

{\rightskip=0pt plus 2em Thus, for \textit{example}, for $D=305$ we have, among others, the period $(17,4,-17)$, $(-17,13,8)$, $(8,11,-23)$, $(-23,12,7)$, $(7,16,-7)$, $(-7,12,23)$, $(23,11,-8)$, $(-8,13,17)$. From it we first choose the form $(7,16,-7)$, and then derive from this the representative form $(7,2,-43)$.\par}

III. When the determinant is a positive square $=kk$, find the reduced form $(A,k,0)$ contained in the proposed class. If $A<k$ or $=k$, take this form itself as the representative. But if $A>k$, take in its place the form $(A-2k,k,0)$, whose first term will be negative but smaller than $k$.

\textit{Example.} In this way all forms with determinant $-235$ are divided into sixteen classes, whose representatives are $(1,0,235)$, $(2,1,118)$, $(4,1,59)$, $(4,-1,59)$, $(5,0,47)$, $(10,5,26)$, $(13,5,20)$, $(13,-5,20)$, and eight others differing from the preceding ones only in the signs of their outer terms: $(-1,0,-235)$, $(-2,1,-118)$ etc.

All forms with determinant $79$ fall into six classes, whose representatives are $(1,0,-79)$, $(3.\,1,-26)$, $(3,-1,-26)$, $(-1,0.\,79)$, $(-3,1,26)$, $-3.\,-1,26)$.

\Needspace{8\baselineskip}
\Article{224}
This classification therefore separates properly equivalent forms completely from the others. Two forms with the same determinant $D$, if they belong to the same class, will be properly equivalent. Any number representable by one can also be represented by the other; and if a number $M$ can be represented by the first form with relatively prime values of the indeterminates\SourceLeaf{230}, the same number can be represented by the other form in the same way, with both representations belonging to the same value of the expression $\sqrt{D}\ (\mathrm{mod.}\ M)$. But if two forms belong to different classes, they will not be properly equivalent. From the representability of a given number by one, we cannot infer its representability by the other. On the contrary, if the number $M$ can be represented by one with relatively prime values of the indeterminates, we are immediately certain that there is no similar representation of the same number by the other form belonging to the same value of the expression $\sqrt{D}(\mathrm{mod.}\ M)$ (see Articles 167, 168).

On the other hand, two forms $F,F'$ from different classes $K,K'$ can indeed be improperly equivalent. In this case, \textit{every} form from one class will be improperly equivalent to \textit{every} form from the other. Every form from $K$ will have its opposite in $K'$, and the classes $K,K'$ themselves will be called \textit{opposite}. Thus in the first example of the preceding article, the third class of forms with determinant $-235$ is opposite to the fourth, and the seventh to the eighth; in the second example, the second class is opposite to the third, and the fifth to the sixth. So, given any two forms from opposite classes, every number $M$ representable by one is also representable by the other. If one representation uses relatively prime values of the indeterminates, the other can do so too, but these two representations belong to opposite values of the expression $\sqrt{D}\ (\mathrm{mod.}\ M)$. --- The rules given above for choosing representative forms are arranged so that opposite classes always receive opposite representative forms.

Finally, there are also classes \textit{opposite to themselves}. If a form and its opposite belong to the same class, it is easy to see that all forms in this class are both properly and improperly equivalent to one another and have their opposites in the same class. Every class containing an ambiguous form has this property; conversely, every class opposite to itself necessarily contains an ambiguous form (Articles 163, 165), and so is called an \textit{ambiguous class}. Thus among the classes of forms with determinant $-235$ there are eight ambiguous ones, represented by $(1,0,235)$, $(2,1,118)$, $(5,0,47)$, $(10,5,26)$, $(-1,0,-235)$, $(-2,1,-118)$, $(-5,0,-47)$, $(-10,5,-26)$; among the classes of forms with determinant $79$ there are two, represented by $(1,0,-79)$, $(-1,0,79)$. --- If the representative forms have been determined by our rules, the ambiguous classes can easily be recognized from them. For a positive\SourceLeaf{231} nonsquare determinant, an ambiguous class certainly receives an ambiguous representative form (Article 194). For a negative determinant, the representative form of an ambiguous class is either itself ambiguous or has equal outer terms (Article 172). Finally, for a positive square determinant, Article 210 makes it easy to decide whether the representative form is improperly equivalent to itself, and thus whether the class it represents is ambiguous.

\Needspace{8\baselineskip}
\Article{225}
We have already shown above (Article 175) that in a form $(a,b,c)$ with negative determinant, the outer terms have the same sign, both as each other and as the outer terms of any equivalent form. If $a,c$ are positive, we call the form $(a,b,c)$ \textit{positive}. We also call the whole class containing $(a,b,c)$, which consists only of positive forms, a \textit{positive class}. Conversely, $(a,b,c)$ will be a \textit{negative form}, contained in a \textit{negative class}, if $a,c$ are negative. Negative numbers cannot be represented by a positive form, nor positive numbers by a negative form. If the form $(a,b,c)$ represents a positive class, the form $(-a,b,-c)$ represents a negative class. It follows that there are as many positive classes as negative ones, and that once the former have been specified, the latter are also known. Thus in investigating forms with negative determinant, it is usually enough to consider the positive classes, whose properties transfer easily to the negative classes.

This distinction applies only to forms with negative determinant. Forms with positive determinant can represent both positive and negative numbers without distinction; indeed, two forms such as $(a,b,c)$, $(-a,b,-c)$ quite often belong to the same class in this case.
